\documentclass[english, parskip=full]{scrartcl}
\usepackage[utf8]{inputenc}  % Kodierung der Datei
\usepackage[english]{babel}  % Sprache auf Deutsch 
\usepackage{color}
\usepackage{graphicx, gensymb}
\usepackage{amsfonts, cancel, amssymb}
\usepackage{amsmath, amsthm, array, esint}
\usepackage{booktabs, multirow}  % schönere Tabellen
\usepackage{caption}  % Anpassung der Captions
\usepackage{scrlayer-scrpage}    % Manipulation des Seitenstils
\pagestyle{scrheadings}  % Stil für die Seite setzen
\automark{section}  % Abschnittsnamen als Seitenbeschriftung 
\setheadsepline{.5pt}    % Kopzeile durch Linie abtrennen
\usepackage[hidelinks]{hyperref}  % Links und weitere PDF-Features
\usepackage{typearea, tikz, pgffor, verbatim}
\usetikzlibrary{calc, arrows.meta,arrows, graphs, quotes}
\usepackage[cache=false, outputdir=build]{minted}
\usemintedstyle{vs}
\usepackage{placeins} % provides \FloatBarrier

\definecolor{bg}{rgb}{0.9,0.9,0.9}
\newcommand\numberthis{\addtocounter{equation}{1}\tag{\theequation}}
\usepackage{svg}
\usepackage{siunitx}
\usepackage{physics}
\usepackage{tensor}
\renewcommand{\bra}{}
\newcommand{\kom}{}
\newcommand{\brak}{}
\renewcommand{\braket}{}
\newcommand{\intx}{}
\newcommand{\intr}{}
\newcommand{\kla}{}
\newcommand{\klam}{}
\newcommand{\klammer}{}
\newcommand{\oper}{}
\newcommand{\tecxt}{}
\newcommand{\Bra}{}
\newcommand{\Ket}{}
\renewcommand{\i}{\text{i}}
\newcommand{\e}{\text{e}}
\newcommand*{\Fourier}{\textsc{Fourier}\ }
\newcommand*{\F}{\mathcal{F}}
\newcommand*{\sinc}{\text{sinc\,}}
\newcommand{\conv}{\mathop{\scalebox{3.5}{\raisebox{-0.38ex}{\makebox[0.5\width][c]{$\ast$}}}}\limits}

\newcommand*{\titel}{01 Lorentz transformations}
\newcommand*{\autor}{Joachim Schwardt}

\hypersetup{pdfauthor={\autor}, pdftitle={\titel}}

%\titlehead{titlehead}
\subject{General Relativity}
\title{\titel}
\author{\autor}
\date{\today}

\begin{document}

%\begin{titlepage}
\maketitle  % Titel setzen
%\tableofcontents  % Inhaltsverzeichnis setzen
%\end{titlepage}

\section*{Task 1: The Bradley problem}
Let earth orbit the sun in the $xy$-plane at the radius $R$ and let $k^\mu$ describe light originating from a star sitting along the $z$-axis at a distance $\ell \gg R$.
Since the star is assumed to be at rest relative to the sun, we have $(k^\mu) = (\omega, 0,k\sin\theta,k\cos\theta)^\top$ in the inertial reference frame (IRF) of the sun.
Of course we also have $\omega = k$ for light.
Here the angle $\theta$ characterizes the finite radius of earth's orbit via $\tan\theta = \frac{R}{\ell}$.
Assuming that the earth is currently moving along the $x$-axis with $\vec{\beta} = \beta\vec{e}_x$, we find the Lorentz transformation to be
\begin{align}
k_\tecxt\text{E}^\mu &= L^\mu_{\ \nu}(\beta\vec{e}_x) k^\nu \equiv \begin{pmatrix}
\gamma & -\beta\gamma & 0 & 0 \\
-\beta\gamma & \gamma & 0 & 0 \\
0 & 0 & 1 & 0 \\
0 & 0 & 0 & 1
\end{pmatrix} \begin{pmatrix}
\omega \\ 0 \\ k\sin\theta \\ k\cos\theta
\end{pmatrix} = \begin{pmatrix}
\gamma\omega   \\ 
-\beta\gamma\omega \\
k\sin\theta \\ k\cos\theta
\end{pmatrix} = \begin{pmatrix}
\omega_\tecxt\text{E} \\ k_{\tecxt\text{E} x}  \\ k_{\tecxt\text{E} y} \\ k_{\tecxt\text{E} z}
\end{pmatrix} .
\end{align}
Therefore the angle $\alpha$ between $\vec{k}_\tecxt\text{E}$ and the $z$-axis is given by
%\begin{align}
%\cos\alpha &= \frac{\vec{k}_\tecxt\text{E}\cdot\vec{e}_z}{k_\tecxt\text{E}} = \frac{k}{\sqrt{ \beta^2 \gamma^2\omega^2 + k^2}} = \frac{1}{\sqrt{\gamma^2\beta^2 + 1}} = \frac{1}{\sqrt{\gamma^2 - 1 + 1}} = \frac{1}{\gamma}.
%\end{align}
%A more useful form of this result is given by
%\begin{align}
%\sin\alpha &= \sqrt{1 - \cos^2\alpha} = \sqrt{1 - \frac{1}{\gamma^2}} = \beta,
%\end{align}
%where we may approximate $\alpha\simeq \beta$ for small velocities.
%In the case of earth, this value is determined by $\beta \approx \frac{\SI{30}{\frac{km}{s}}}{\SI{3e5}{\frac{km}{s}}} = 10^{-4}$.
%This is equivalent to $\alpha\approx \frac{360^\circ}{2\pi}\cdot 10^{-4} \approx 20.6''$.
%%% TANGENT
%\begin{align}
%\tan\alpha &= \frac{\sqrt{k_x^2 + k_y^2}}{k_z} = \frac{\sqrt{\gamma^2\beta^2\omega^2 + k^2\sin^2\theta}}{k\cos\theta} = \sqrt{\frac{\gamma^2\beta^2+1}{\cos^2\theta} - 1} = \\
%&= \sqrt{\gamma^2 \kla\left( 1 + \tan^2\theta \right) - 1} = \sqrt{\gamma^2\beta^2 + \gamma^2\tan^2\theta} = \\
%&= \gamma\sqrt{\beta^2 + \frac{R^2}{L^2}}.
%\end{align}
%%% SINUS
\begin{align}
\sin\alpha &= \frac{\sqrt{k_x^2 + k_y^2}}{\sqrt{k_x^2 + k_y^2 + k_z^2}} = \frac{\gamma^2\beta^2\omega^2 + k^2\sin^2\theta}{\sqrt{\gamma^2\beta^2\omega^2 + k^2}} = \\
&= \frac{\sqrt{\gamma^2\beta^2 + \sin^2\theta}}{\gamma} = \sqrt{\beta^2 + (1-\beta^2)\sin^2\theta}.
\end{align}
where we may approximate $\alpha\simeq \beta$ for small velocities and $\theta \approx 0$ for a distant source.
In the case of earth we find $\beta \approx \frac{\SI{30}{\frac{km}{s}}}{\SI{3e5}{\frac{km}{s}}} = 10^{-4}$.
This is equivalent to $\alpha\approx \frac{360^\circ}{2\pi}\cdot 10^{-4} \approx 20.6''$.

\section*{Task 2: Four-vector products}
We know that $(u^\mu) = (\gamma,\gamma\vec{\beta})^\top$,  $p^\mu = mu^\mu$, $(k^\mu) = (\omega, \vec{k})^\top$ and $(x^\mu) = (t,\vec{x})^\top$.
Therefore
\begin{align}
u_\mu u^\mu &= \gamma^2 - \gamma^2\beta^2 = 1,\\
p_\mu p^\mu &= m^2 u_\mu u^\mu = m^2,\\
x_\mu k^\mu &= \omega t - \vec{k}\cdot\vec{x}.
\end{align}
Let $a^\mu,b^\mu$ be arbitrary four-vectors and let the Lorentz-transformation be
\begin{align}
(\Lambda^\mu_{\ \nu}) &= \begin{pmatrix}
\gamma & \beta\gamma & 0 & 0 \\
\beta\gamma & \gamma & 0 & 0 \\
0 & 0 & 1 & 0 \\
0 & 0 & 0 & 1
\end{pmatrix},
\end{align}
which transforms the four-vectors into $\tilde{v}^\mu = \Lambda^\mu_{\ \nu} v^\nu$ (for $v\in\{a,b\}$).
Then we find
\begin{align}
\tilde{a}_\mu \tilde{b}^\mu &= \tilde{a}^\nu g_{\nu\mu} \tilde{b}^\mu = \Lambda^\nu_{\ \alpha}a^\alpha g_{\nu\mu} \Lambda^\mu_{\ \beta} b^\beta = a^\alpha \Lambda_{\alpha}^{\ \nu} g_{\nu\mu} \Lambda^\mu_{\ \beta} b^\beta.
\end{align}
Note that
\begin{align}
\Lambda_{\alpha}^{\ \nu} g_{\nu\mu} \Lambda^\mu_{\ \beta} &\equiv \begin{pmatrix}
\gamma & \beta\gamma & 0 & 0 \\
\beta\gamma & \gamma & 0 & 0 \\
0 & 0 & 1 & 0 \\
0 & 0 & 0 & 1
\end{pmatrix} \begin{pmatrix}
1 & 0 & 0 & 0 \\
0 & -1 & 0 & 0 \\
0 & 0 & -1 & 0 \\
0 & 0 & 0 & -1
\end{pmatrix} \begin{pmatrix}
\gamma & \beta\gamma & 0 & 0 \\
\beta\gamma & \gamma & 0 & 0 \\
0 & 0 & 1 & 0 \\
0 & 0 & 0 & 1
\end{pmatrix} \\ 
&= \begin{pmatrix}
\gamma & \beta\gamma & 0 & 0 \\
\beta\gamma & \gamma & 0 & 0 \\
0 & 0 & 1 & 0 \\
0 & 0 & 0 & 1
\end{pmatrix} \begin{pmatrix}
\gamma & \beta\gamma & 0 & 0 \\
-\beta\gamma & -\gamma & 0 & 0 \\
0 & 0 & -1 & 0 \\
0 & 0 & 0 & -1
\end{pmatrix} = \\
&= \begin{pmatrix}
\gamma^2 - \beta^2\gamma^2 & \gamma\beta\gamma + \beta\gamma(-\gamma) & 0 & 0 \\
\beta\gamma^2 + \gamma(-\beta\gamma) & \beta\gamma\beta\gamma + \gamma(-\gamma) & 0 & 0 \\
0 & 0 & -1 & 0 \\
0 & 0 & 0 & -1
\end{pmatrix} = \\
&= \begin{pmatrix}
1 & 0 & 0 & 0 \\
0 & -1 & 0 & 0 \\
0 & 0 & -1 & 0 \\
0 & 0 & 0 & -1
\end{pmatrix} \equiv g_{\alpha\beta}.
\end{align}
Using this identity, the previous product simplifies to
\begin{align}
\tilde{a}_\mu \tilde{b}^\mu &= a^\alpha g_{\alpha\beta} b^\beta = a_\alpha b^\alpha.
\end{align}
Thus four-vector products are always invariant under Lorentz transformation.

\section*{Task 3: Electrodynamics}
The field strength tensor is defined as $F_{\mu\nu} := \partial_\mu A_\nu - \partial_\mu A_\nu$ and is hence fully anti-symmetric.
Thus we only need to consider the components $\mu>\nu$.
For $\nu = 0$ and $j\in\{1,2,3\}$ we find
\begin{align}
F_{j 0} &= \partial_jA_0 - \partial_0 A_j \equiv \nabla \phi + \partial_t\vec{A} = -\vec{E}.
\end{align}
For $\nu > 0$ and $j\neq k$ we find \footnote{Note that $\sum_a\epsilon_{abc}\epsilon_{ab'c'} = \delta_{bb'}\delta_{cc'} - \delta_{bc'}\delta_{cb'}$.}
\begin{align}
F_{jk} &= \partial_j(-A_k) - \partial_k (-A_j) =  \sum_{bc} \kla\left( \delta_{jb}\delta_{kc} - \delta_{jc}\delta_{kb} \right)  \partial_{b} (-A_{c}) = \\
&= -\sum_{abc}\epsilon_{ajk}\epsilon_{abc} \partial_{b}A_{c} = -\sum_a \epsilon_{ajk} B_a.
\end{align}
Therefore we may write
\begin{align}
F_{\mu\nu} &= \begin{pmatrix}
0 & \vec{E} \\ -\vec{E} & -(\epsilon_{jkl}B_l)_{jk}
\end{pmatrix} &&\tecxt\text{and} & F^{\mu\nu} &= \begin{pmatrix}
0 & -\vec{E} \\ \vec{E} & -(\epsilon_{jkl}B_l)_{jk}
\end{pmatrix}.
\end{align}
In this form we can elegantly compute the contraction
\begin{align}
F^{\mu\nu}F_{\mu\nu} &= -2\vec{E}^2 + \sum_{jkll'}\epsilon_{jkl}\epsilon_{jkl'}B_lB_{l'} = \\
&= -2\vec{E}^2 + \sum_{kll'} \kla\left( \delta_{ll'} - \delta_{kl'}\delta_{kl} \right)B_lB_{l'} = 2\kla\left( \vec{B}^2 - \vec{E}^2 \right).
\end{align}
We may also express the equations $\partial_\mu F^{\mu\nu}=j^\nu$ and $\partial_\mu F_{\nu\rho} + \partial_\nu F_{\rho\mu} + \partial_\rho F_{\mu\nu} = 0$ in terms of $\vec{E}$ and $\vec{B}$.
For $\nu=0$ the first equation is
\begin{align}
\partial_\mu F^{\mu 0} &= \nabla \cdot \vec{E} = \rho.
\end{align}
For $\nu > 0$ the first equation is
\begin{align}
\partial_\mu F^{\mu j} &= \partial_t F^{0j} + \partial_k F^{kj} = \partial_t \kla\left( -\vec{E} \right) + \partial_k \kla\left( -\epsilon_{kjl}B_l \right) = \\
&= \nabla\times\vec{B} - \partial_t\vec{E} = \vec{j}.
\end{align}
For $\rho=0$ and $\mu>\nu > 0$ the second equation is
\begin{align}
0 &= \partial_\mu F_{\nu 0} - \partial_\nu F_{\mu 0} + \partial_0 F_{\mu \nu} \equiv -\partial_\mu E_\nu + \partial_\nu E_\mu - \partial_t \epsilon_{\mu\nu k}B_k = \\
&= -\epsilon_{\mu\nu k} (\nabla\times \vec{E})_k - \epsilon_{\mu\nu k}\partial_t B_k,\\
\Rightarrow\ \ \ 0 &= \nabla\times\vec{E} + \partial_t\vec{B}.
\end{align}
For $\rho>\mu>\nu> 0$ the second equation is
\begin{align}
0 &= \partial_\mu \kla\left( -\epsilon_{\nu\rho k}B_k \right) + \partial_\nu \kla\left( -\epsilon_{\rho\mu k}B_k \right) + \partial_\rho \kla\left( -\epsilon_{\mu\nu k} B_k \right) = \\
&= -\kla\left( \partial_2\epsilon_{13k} + \partial_1\epsilon_{32k} + \partial_3\epsilon_{21k} \right) B_k = -\partial_k B_k,\\
\Rightarrow\ \ \ 0 &= \nabla\cdot\vec{B}.
\end{align}

%\begin{align}
%0 &= \partial_\mu F_{\nu\rho} + \partial_\nu F_{\rho\mu} + \partial_\rho F_{\mu\nu} = \\
%&= \partial_\mu \partial_\nu A_\rho - \partial_\mu \partial_\rho A_\nu + \partial_\nu \partial_\rho A_\mu - \partial_\nu \partial_\mu A_\rho + \partial_\rho \partial_\mu A_\nu - \partial_\rho \partial_\nu A_\mu.
%\end{align}
%For $\rho=0$ this reduces to
%\begin{align}
%0 &= \partial_t\partial_\nu A_\mu - \partial_t \partial_
%\end{align}


%\begin{thebibliography}{9}
%\bibitem{example} description, \url{https://www.exmapleurl}, day.\,month~2021.
%\end{thebibliography}

\end{document}